\documentclass[11pt,letterpaper]{article}
\usepackage[T1]{fontenc}
\usepackage{lmodern,amsmath,amssymb,amsthm,mathtools,microtype}
\usepackage[margin=1in,top=.9in,bottom=.92in]{geometry}
\usepackage{xcolor,booktabs,fancyhdr,enumitem}
\usepackage[colorlinks=true,linkcolor=teal!50!black,urlcolor=teal!60!black,citecolor=teal!55!black]{hyperref}
\hypersetup{pdftitle={Dirichlet characters from a single Li heat kernel},pdfauthor={Research note}}
\definecolor{accent}{HTML}{125A61}
\newcommand{\R}{\mathbb R}\newcommand{\T}{\mathbb T}\newcommand{\Z}{\mathbb Z}
\newcommand{\C}{\mathbb C}\newcommand{\N}{\mathbb N}\newcommand{\dd}{\,\mathrm d}
\newcommand{\I}{\mathrm i}\newcommand{\e}{\mathrm e}
\newcommand{\wh}{\widehat}\newcommand{\one}{\mathbf 1}
\DeclareMathOperator{\supp}{supp}\DeclareMathOperator{\dist}{dist}
\theoremstyle{plain}
\newtheorem{theorem}{Theorem}[section]\newtheorem{lemma}[theorem]{Lemma}
\newtheorem{proposition}[theorem]{Proposition}\newtheorem{corollary}[theorem]{Corollary}
\theoremstyle{definition}\newtheorem{example}[theorem]{Example}
\newcommand{\artbase}{https://github.com/VladimirReshetnikov/ProveIt/blob/1a50d416ba30b91595db12d7437f0328a32b0468/Analysis/GeneralizedRiemannHypothesis}
\newcommand{\leanbase}{https://github.com/lsernau/grh-lean/blob/72fead102f663d5c7a407964a13dc0cf9d977207/grh/src}
\setlist{nosep,topsep=4pt}
\setlength{\headheight}{15pt}
\pagestyle{fancy}\fancyhf{}\lhead{\small\textsc{Dirichlet characters and Li heat kernels}}\rhead{\small\thepage}
\renewcommand{\headrulewidth}{.3pt}
\raggedbottom
\setlength{\emergencystretch}{2em}
\begin{document}
\thispagestyle{empty}
\noindent{\sffamily\small\color{accent} ANALYTIC NUMBER THEORY \hfill 27 NOVEMBER 2026}
\vspace{8mm}
\begin{center}
{\LARGE\bfseries Dirichlet characters from a single Li heat kernel\par}
\vspace{3mm}
{\large Signed zeros and a sharp small-time law on the circle\par}
\end{center}
\vspace{3mm}
\begin{abstract}\noindent
For a primitive Dirichlet character $\chi$, exponentiating its full complex Li coefficients gives a positive convolution semigroup on the circle. One exact density at one unknown positive time determines $\chi$, even with an unknown angular origin. Magnitudes of its Fourier coefficients determine the character up to conjugation. Both assertions persist under relative agreement of the Fourier tails outside a set of density zero. At short times the height of the density at an individual zero angle, relative to its height at the origin, reads off the signed zero's multiplicity. At a shortest sum of $k$ angles the corresponding coefficient is the number of ordered representations divided by $k!$. At almost every other point the log of the reciprocal relative height is asymptotic to $L\log L/\log 2$, where $L=\log(1/t)$. The same constant governs the minimum over the circle. For products of Cayley factors, angular precision $10^{-D}$ costs asymptotically $\log_2 D$ zero factors for almost every target, with a matching worst-case count. An interval computation for a conjugate pair modulo five certifies thirteen successive asymmetric examples. The arithmetic input is the formal zero-line theorem for Dirichlet $L$-functions.
\end{abstract}

\section{The inverse problem}
Write $\T=\R/(2\pi\Z)$, with ordinary angular Lebesgue measure $\dd y$. Under the Dirichlet zero-line theorem every nontrivial zero has the form $\rho=1/2+\I\gamma$. Its Cayley image and inverse are
\begin{equation}\label{cayley}
 z_\rho=1-\rho^{-1}=\e^{\I\theta_\gamma},\qquad
 \theta_\gamma=2\arctan(1/(2\gamma)),\qquad
 \gamma=\tfrac12\cot(\theta_\gamma/2).
\end{equation}
We set $\theta_0=\pi$. Ordinates and angles are signed. The point $0$ on the circle is a limit of angles and corresponds to no finite zero. We use natural logarithms except for $\log_2 u=(\log u)/(\log 2)$.

Use the zero-sum convention for the Li coefficients of a primitive character of conductor $q$, including the trivial character of conductor one for $\zeta$:
\begin{equation}\label{li}
 \lambda_\chi(n)=\lim_{T\to\infty}\sum_{|\gamma|\le T}
                (1-z_\rho^n),\qquad n\in\Z,\qquad
 p_{\chi,t}(y)=\frac1{2\pi}\sum_{n\in\Z}
                  \exp(-t\lambda_\chi(n)-\I n y),\quad t>0.
\end{equation}
Zeros are repeated by multiplicity. Symmetric truncation matters for the imaginary part. In this convention $\lambda_\chi(-n)=\overline{\lambda_\chi(n)}$. Using instead the conventional logarithmic derivatives at $s=1$ reverses the index in the zero-sum formula for a complex character \cite{lagarias}. A reflection of the angular coordinate accounts for that convention.

Three features of \eqref{li} are worth distinguishing. Taking the logarithm of each Fourier coefficient loses integer multiples of $2\pi\I$. Nevertheless the entire sequence determines the character at a single time (Section~\ref{inverse}). As $t$ decreases, values of the density at specified points determine the individual signed multiplicities and their additive relations (Section~\ref{layers}). At a typical angle, finite sums never land exactly on target. The cost of successively improving an approximation determines the sharp decay rate (Section~\ref{typical}). The constants in small-time estimates may depend on the fixed character.

\paragraph{The repository inputs.}
The public statements in \texttt{grh/src/Statements.lean} use Mathlib's Dirichlet characters and $L$-function. The declaration \texttt{LFunction\_eq\_zero\_re} in \texttt{Main.lean} places all open-strip zeros on the critical line for every positive modulus. Its primitive specialization supplies \eqref{cayley}. Sernau and collaborators' rewriting of Alp\"oge's proof derives zero exclusion from a one-sided bound on a smoothed prime sum, obtained via lattice covolumes. Sutherland records that both editions compiled, that every module was replayed, and that the five public statements passed independent checkers under the three standard Lean axioms \cite{lean,mit}.

The ProveIt article \cite{article} merges 103 reports. Its section on Li coefficients uses the symmetric measure with half a jump at each of $\theta_\gamma$ and $-\theta_\gamma$ for each zero. It obtains a central Cauchy profile and higher logarithmic corrections. We retain a full unit jump at $\theta_\gamma$. This changes what can be distinguished at a fixed time and what the small-time density measures off the origin. Nearby ProveIt reports already use finite lists of low ordinates to distinguish characters \cite{lowzeros}. Our data are the Li heat density or the Li coefficients. The arguments below are mathematical proofs outside the upstream Lean development.

\section{The oriented semigroup and the scale near the origin}\label{construction}
Set $\nu=\sum_\gamma\delta_{\theta_\gamma}$ with multiplicities and use representatives in $(-\pi,\pi]$. The one-sided zero counts, for either sign of $\gamma$, satisfy
\begin{equation}\label{rv}
 N_\chi^\pm(T)=\frac{T}{2\pi}\log\frac{qT}{2\pi\e}+O_\chi(\log(2+T)).
\end{equation}
The zero-sum formula and this one-sided estimate are discussed in \cite[Theorems 1.1 and 2.1]{lagarias}. In particular $\int\theta^2\dd\nu<\infty$. The nonzero support is discrete off $0$, with at least one positive angle arbitrarily close to $0$.

\begin{proposition}\label{process}
The function $p_{\chi,t}$ is an entire periodic function, strictly positive on $\T$, of integral one. Its Fourier coefficients in the convention $\wh p(n)=\int_\T p(y)\e^{\I ny}\dd y$ are $\e^{-t\lambda_\chi(n)}$. The laws form a convolution semigroup with jump measure $\nu$.
\end{proposition}
\begin{proof}
For a finite cutoff the law of the sum of independent Poisson jumps of unit rate at the listed angles has the stated characteristic function with the truncated exponent. Its pointwise limit is positive definite on $\Z$ and equals one at $n=0$, hence defines a probability on $\T$. More explicitly the compensated real-line exponent, taken with a minus sign in the characteristic function, is
\[
 \int(1-\e^{\I u\theta}+\I u\theta)\dd\nu(\theta)-\I ub,\qquad
 b=\lim_{T\to\infty}\sum_{|\gamma|\le T}\theta_\gamma.
\]
The limit exists since $\sum^*\sin\theta_\gamma=-\Im\lambda_\chi(1)$ and $\sum|\theta_\gamma-\sin\theta_\gamma|<\infty$. At integers the exponent equals \eqref{li}. This construction also specifies the drift after a cutoff of large jumps.

Lagarias's degree-one GRH asymptotic gives, as $n\to+\infty$,
\begin{equation}\label{asym}
 \lambda_\chi(n)=\tfrac12 n\log n+C_\chi n+O_\chi(\sqrt n\log(2+n)),
 \qquad C_\chi\in\R .
\end{equation}
Thus the series \eqref{li} converges absolutely on each horizontal strip and is the continuous density of the probability constructed above. The semigroup property is Fourier multiplication. A nonzero real-analytic nonnegative density is positive almost everywhere. Convolving the half-time density with itself gives strict positivity at every point.
\end{proof}

Let $\gamma_E$ be Euler's constant, $W$ the positive inverse of $w\mapsto w\e^w$, and
\begin{equation}\label{a}
 w=W\!\left(\frac{q\e^{\gamma_E-1}}{\pi t}\right),\qquad
 a=a(t)=\tfrac12 t w\sim\tfrac12 t\log(1/t).
\end{equation}
For bounded $x$ the central profile is
\begin{equation}\label{centre}
 a\,p_{\chi,t}(ax)\ \longrightarrow\ \frac{1}{\pi(1+x^2)}
 \quad\text{locally uniformly in }x,\qquad
 \frac1{2\pi}\sum_{n\in\Z}\e^{-t\Re\lambda_\chi(n)}
           \sim p_{\chi,t}(0)\sim\frac1{\pi a}.
\end{equation}
Indeed the profile for $\Re\lambda_\chi$ is Proposition \texttt{li:prop:trace} of \cite{article}. Replacing it by $\lambda_\chi$ changes the scaled density by at most
\[
 \frac{at}{2\pi}\sum_n|\Im\lambda_\chi(n)|\e^{-t\Re\lambda_\chi(n)}
       =O_\chi(\sqrt a).
\]
To see the bound, on $|n|\ge a^{-1/2}$ the real exponential is at most $\e^{-c a|n|}$ by \eqref{asym}; sum $\sqrt n\log(2+n)$ against this majorant and use $t\log(1/a)\asymp a$. The smaller modes contribute $O(\sqrt a)$ as well by direct summation. Multiplying the summands by any fixed power of $a|n|$ gives the same order of error. Thus even the higher central logarithmic corrections of \cite{article} survive the orientation. For \eqref{centre} alone one can directly use \eqref{asym} and Riemann sums at mesh $a$: $t\lambda_\chi(\lfloor s/a\rfloor)\to |s|$, giving the inverse Fourier transform of $\e^{-|s|}$.

\section{A single time determines the character}\label{inverse}
Convergence in natural density of a sequence $u_n$ to $u$ means that for each $\varepsilon>0$ the proportion of $1\le n\le N$ with $|u_n-u|>\varepsilon$ tends to zero. Write $F_{\chi,t}(n)=\exp(-t\lambda_\chi(n))$. In particular none of the denominators below vanish.

\begin{theorem}[One-time rigidity]\label{single}
Let $\chi,\psi$ be primitive Dirichlet characters, $s,t>0$, and $\alpha\in\T$.
\begin{enumerate}[label=\textup{(\roman*)}]
\item If $F_{\chi,t}(n)/(\e^{\I n\alpha}F_{\psi,s}(n))\to1$ in natural density on positive integers, then $s=t$, $\chi=\psi$ and $\alpha=0$ in $\T$.
\item If $|F_{\chi,t}(n)|/|F_{\psi,s}(n)|\to1$ in natural density, then $s=t$ and $\psi\in\{\chi,\bar\chi\}$. Both possibilities give the same magnitudes.
\end{enumerate}
One exact time slice, even restricted to a nonempty open arc and known up to rotation, identifies its time and character. If the angular orientation is also unknown the ambiguity is precisely conjugation.
\end{theorem}

\begin{lemma}[Finite zero discrepancies]\label{finite}
Suppose two finite products of primitive Dirichlet $L$-functions have nontrivial zero multisets differing at only finitely many occurrences, after cancellation. Then their uncompleted products are identical.
\end{lemma}
\begin{proof}
Use the entire order-one completions, multiplying the completed zeta factor by $s(s-1)$ when it occurs. Their zeros are exactly the nontrivial zeros. Hadamard factorization gives a quotient $\exp(A+Bs)R(s)$ with $R$ rational. Write $D(s)$ for the quotient of the uncompleted products. The completed factors differ from the uncompleted ones by powers of constants, exponentials of linear polynomials, possible polynomial factors, and factors $\Gamma((s+\kappa)/2)$ with $\kappa=0$ or $1$. The gamma recurrence therefore makes $D(s+2)/D(s)$ a constant times a rational function. On the positive real axis $D(\sigma)=1+O(2^{-\sigma})$. A rational function approaching one with exponential accuracy is identically one, by expansion at infinity. Thus $D(s+2)=D(s)$, and shifting arbitrarily far right gives $D=1$ on $\Re s>1$. Meromorphic continuation finishes the proof.
\end{proof}

\begin{proof}[Proof of Theorem~\ref{single}]
An unbounded subsequence of convergence and the leading real term of \eqref{asym} force $s=t$ in either case. In fact $t=\lim_{n\to\infty}-2\log|F_{\chi,t}(n)|/(n\log n)$. Put $\Delta^2 f(n)=f(n+1)-2f(n)+f(n-1)$ and
\begin{equation}\label{curv}
 D_n=\Delta^2(\lambda_\chi-\lambda_\psi)(n)=\int_\T\e^{\I n\theta}\dd\eta(\theta),
 \qquad \eta=2(1-\cos\theta)(\nu_\chi-\nu_\psi).
\end{equation}
The weight makes the two weighted measures finite. The sequence $D_n$ is a uniform limit of trigonometric polynomials on $\Z$ because the measures are purely atomic. It is therefore uniformly almost periodic. For such a sequence, shifts changing every term by less than a specified tolerance have bounded gaps. This follows for finitely many frequencies from compactness of the closure of their joint rotations and passes to uniform limits. Continuous functions of bounded almost periodic sequences are almost periodic as well.

In (i), take the multiplicative second difference of the proposed ratios. The rotation factor disappears, giving $\exp(-tD_n)\to1$ in density, since shifting a zero-density exceptional set by one preserves density zero. An almost periodic sequence which converges to a constant in density equals that constant everywhere: a single deviation recurs within half its size on a set of shifts with bounded gaps and hence positive lower density. Consequently $D_n$ takes values in $(2\pi\I/t)\Z$. An almost period with tolerance smaller than $2\pi/t$ is an exact period. Choose a positive such period $r$. Uniqueness of Fourier coefficients of finite measures now gives
\[
                 (\e^{\I r\theta}-1)\eta=0 .
\]
The measure $\eta$ is supported on the finitely many $r$-th roots of unity. The weight in \eqref{curv} vanishes only at $0$, where there is no zero angle. The two original zero multisets thus have at most finitely many discrepant occurrences. Lemma~\ref{finite}, Dirichlet-series uniqueness and primitivity give $\chi=\psi$. The remaining ratio $\e^{-\I n\alpha}$ converges to one in density. Its Ces\`aro mean must be one, which forces $\e^{\I\alpha}=1$.

In (ii) the real logarithm is single valued, and $\Re(\lambda_\chi-\lambda_\psi)(n)\to0$ in density. Its second differences are almost periodic and thus identically zero. By \eqref{curv} the symmetrizations of the two angle measures agree off zero. Hence $L(s,\chi)L(s,\bar\chi)$ and $L(s,\psi)L(s,\bar\psi)$ have identical zero multisets and, by Lemma~\ref{finite}, coincide. At unramified primes their first Dirichlet coefficients give $\chi(p)+\bar\chi(p)=\psi(p)+\bar\psi(p)$. Dirichlet's theorem supplies a prime in every reduced class of the joint modulus. Character orthogonality then gives equality of the pairs of induced characters with multiplicities, and uniqueness of primitive induction gives (ii). Conversely conjugation reflects $\gamma$ and conjugates $\lambda(n)$, preserving its real part.

Equality on an open arc extends to the whole circle by analyticity. Reflection sends $p_{\chi,t}(y)$ to $p_{\bar\chi,t}(y)$, completing the assertions about coordinates.
\end{proof}

The relative accuracy in Theorem~\ref{single} is essential to its interpretation. By \eqref{asym} the high Fourier coefficients themselves are extremely small. Also a general circle semigroup can lose its generator at a single time \cite{hasebe}. For example, take Poisson jumps of rate $\pi/t$ at $\pi/2$ or at $-\pi/2$. The ratio of their time-$t$ Fourier coefficients is $\exp(2\pi\I\sin(n\pi/2))=1$ for every integer $n$. In the arithmetic problem finite zero rigidity removes exactly this type of obstruction.

\begin{corollary}[A bounded, rotation-invariant observation]\label{bounded}
The sequence $C_{\chi,t}(n)=F_{\chi,t}(n+1)F_{\chi,t}(n-1)/F_{\chi,t}(n)^2$ alone determines $t$ and $\chi$. Its modulus alone determines $t$ and the conjugate pair. These conclusions hold if the ratio of two such observations tends to one in natural density.
\end{corollary}
\begin{proof}
Both $C_{\chi,t}$ and its reciprocal are bounded almost periodic sequences by \eqref{curv} applied to one measure. In the complex case write $D=\Delta^2(t\lambda_\chi-s\lambda_\psi)$. The proof of Theorem~\ref{single} makes $\exp(-D_n)=1$, then makes $D$ periodic. Outside finitely many angles the measures $t\nu_\chi$ and $s\nu_\psi$ agree. In the modulus case the same argument gives $\Re D_n=0$ and agreement of their symmetrizations everywhere. Counting the angles corresponding to $|\gamma|\le T$ yields $tN_\chi(T)-sN_\psi(T)=O(1)$. The common leading term $T\log T/\pi$ in \eqref{rv} forces $t=s$. The remaining assertions follow as before.
\end{proof}

\section{Heights at sums of zero angles}\label{layers}
Let $S_0=\{0\}$ and let $S_j$ consist of sums modulo $2\pi$ of at most $j$ zero angles. These are compact, since the atoms together with $0$ form a compact set. For $y\in S_k\setminus S_{k-1}$ write $M_k(y)$ for the number of ordered $k$-tuples of zero occurrences with angle sum $y$. A zero of multiplicity $m$ gives $m$ choices at each use, and occurrences may be selected repeatedly.

\begin{theorem}[A multiplicity formula at every finite depth]\label{echo}
Fix $y\in S_k\setminus S_{k-1}$, $k\ge1$. The integer $M_k(y)$ is finite and positive. Uniformly on each bounded set of real $x$,
\begin{align}
 \frac{p_t(y+a x)}{p_t(a x)}
   &=\frac{M_k(y)}{k!}\,t^k+O_{\chi,y,x_0}(t^{k+1})\qquad(|x|\le x_0),\label{echoratio}\\
 \frac{p_t(y+a x)}{t^kp_t(0)}
   &\longrightarrow \frac{M_k(y)}{k!(1+x^2)} .\label{echoprof}
\end{align}
If $y\notin\bigcup_jS_j$, the ratio in \eqref{echoratio} is $O_{\chi,y,x_0,N}(t^N)$ for every integer $N$.
\end{theorem}
\begin{proof}
Let $d=\dist(y,S_{k-1})>0$, for arc distance on the circle. If a sum of $k$ angles is within $d/2$ of $y$, deleting any summand leaves a point of $S_{k-1}$, so every summand has norm at least $d/2$. Only finitely many such angles and occurrences exist. Thus $M_k(y)$ is finite and $y$ has positive distance from $S_k\setminus\{y\}$. Choose $h>0$ smaller than one quarter of both distances.

For a fixed small $\eta>0$, let $\mu$ be the finite part of $\nu$ on angles with norm greater than $\eta$, and let $\ell=\mu(\T)$. Remove the corresponding compound Poisson jumps and write $g_t$ for the circle density of the remainder. Removing them changes the Fourier exponents by at most $2\ell$; \eqref{centre} gives
\begin{equation}\label{smallbound}
 \|g_t\|_\infty\ll_\eta a^{-1},\qquad
 \|g_t-p_t\|_\infty\ll_\eta t/a .
\end{equation}
At a fixed distance from zero we can obtain any prescribed power of $t$ by choosing $\eta$ small enough. Indeed, lift the residual process to $Y_v$ on the line. With $b_\eta=b-\int_{|\theta|>\eta}\theta\dd\nu$ and $V=\int\theta^2\dd\nu$, the compensated exponential formula gives
\[
 \mathbb E\e^{\pm sY_v}
 \le\exp\!\left(vs|b_\eta|+\tfrac12 v s^2\e^{s\eta}V\right).
\]
Taking $s=\log(1/t)/(2\eta)$ bounds both moments uniformly for $v=t/2$. Therefore $\Pr(|Y_{t/2}|\ge h/2)\ll t^{h/(4\eta)}$. In $g_{t/2}*g_{t/2}(z)$ for $\dist(z,0)\ge h$ at least one argument has norm at least $h/2$. Bound the other density by \eqref{smallbound} and use $a(t/2)\asymp a(t)$. This gives $\sup_{\dist(z,0)\ge h}g_t(z)\ll_\eta a^{-1}t^{h/(4\eta)}$.

Now choose $\eta<d/2$ and $h/(4\eta)>k+1$. Expand the full density as
\[
 p_t=\e^{-\ell t}\sum_{j\ge0}\frac{t^j}{j!}\,\mu^{*j}*g_t,
 \qquad\mu^{*0}=\delta_0.
\]
For $j<k$ all centers stay away from $y$, as do the nonmatching centers at $j=k$. The preceding bound applies uniformly also at $y+ax$ for bounded $x$ and small $t$. At the matching center the coefficient is exactly $M_k(y)$. Terms with $j>k$ cost $O_\eta(t^{k+1}/a)$ by \eqref{smallbound} and the exponential series. Replacing $\e^{-\ell t}g_t(ax)$ by $p_t(ax)$ costs $O_\eta(t/a)$. Since $p_t(ax)\gg_{x_0}a^{-1}$, this proves \eqref{echoratio}, and \eqref{centre} gives \eqref{echoprof}. At a point outside all finite sums, use $\dist(y,S_N)>0$ and the same argument to rule out all terms of orders $0$ through $N$.
\end{proof}

For every nonzero $y$ on the circle, in particular,
\begin{equation}\label{readzero}
 \lim_{t\downarrow0}\frac{p_{\chi,t}(y)}{t\,p_{\chi,t}(0)}
 =\operatorname{ord}_{\,s=1/2+(\I/2)\cot(y/2)}L(s,\chi).
\end{equation}
This includes a central zero at $y=\pi$. The limit of $\log(p_t(y)/p_t(0))/\log t$ is the least number of zero factors whose Cayley images multiply to $\e^{\I y}$, with value $+\infty$ if there are none. The finite-depth points form a countable dense subset. Indeed a sufficiently small positive angle and its successive multiples give an arbitrarily fine net. Every integer depth occurs: equality of two successive $S_j$ would stabilize all subsequent sums, whereas a compact countable set cannot be dense in the circle.

\begin{corollary}[Finite Li data and an error margin]\label{finite-li}
Fix $k\ge1$ and $c>2k$. Given the coefficients with $|n|\le N$, use $t=c/N$ and truncate \eqref{li} to these frequencies, obtaining $P_N$. For $y\notin S_{k-1}$ the limit of $k!P_N(y)/(t^kP_N(0))$ is $M_k(y)$, understood as zero off $S_k$. With Hermitian input errors $\max_{|n|\le N}|\widetilde\lambda(n)-\lambda(n)|\le\epsilon_N$, rounding to the nearest integer still succeeds eventually at that $y$ if
\[
                 \limsup_{N\to\infty}\frac{\epsilon_N}{t^{k-1}}<\frac1{2k!}.
\]
In particular the first-layer formula \eqref{readzero} tolerates any uniform coefficient error strictly below $1/2$.
\end{corollary}
\begin{proof}
Choose $b<1/2$ with $bc>k$. For large $N$, \eqref{asym} bounds the omitted absolute Fourier tail by
$O(\sum_{n>N}\exp(-bt n\log N))=O(a^{-1}N^{-bc})=o(t^k/a)$.
An input perturbation changes the polynomial uniformly by at most
$t\epsilon_N\e^{t\epsilon_N}(2\pi)^{-1}\sum_n\e^{-t\Re\lambda(n)}
=(1+o(1))t\epsilon_Np_t(0)$ under the stated error condition. The relative error of the denominator is $o(1)$. Apply Theorem~\ref{echo} to the true ratio, and multiply the absolute numerator error by $k!/(t^k p_t(0))$. The limiting error is strictly below $1/2$. Hermitian inputs ensure a real statistic. This is pointwise eventual recovery; its starting $N$ can depend on the queried angle.
\end{proof}
\section{A sharp law at almost every point}\label{typical}
Fix a primitive character and write $p_t$ for its oriented circle kernel,
$a=a(t)$ for its central spatial scale, and $L=\log(1/t)$.
All constants in this section may depend on the character.
\begin{theorem}[Typical density and the minimum]\label{sharp-heat}For $t$ sufficiently small one has
\begin{equation}\label{eq:uniformfloor}
 \inf_{y\in\mathbb T}\frac{p_t(y)}{p_t(0)}
 \ \ge\ \exp\!\left(-\frac{L\log L}{\log 2}-O(L)\right).
\end{equation}
For Lebesgue-almost every fixed $y\in\mathbb T$,
\begin{equation}\label{eq:typicallaw}
 \lim_{t\downarrow0}
 \frac{-\log(p_t(y)/p_t(0))}{L\log L}=\frac1{\log 2}.
\end{equation}
The same limit holds with $p_t(y)$ replaced by $\min_y p_t(y)$.\end{theorem}
The first estimate gives one inequality for this last assertion, and any
one point satisfying \eqref{eq:typicallaw} gives the other.

\paragraph{Uniform lower bound.}
Write $N_+(T)$ for the count of positive ordinates up to $T$. Its main
term is $(T/(2\pi))\log(qT/(2\pi e))$ with error $O_\chi(\log(T+2))$.
Taking a sufficiently large fixed $D_\chi$ in its first difference
shows every $[T,T+D_\chi]$ contains a positive ordinate for large $T$.
As $d\theta/dT=-1/(T^2+1/4)$, there are constants $r_0,C>0$ such
that every $0<r\le r_0$ admits a positive jump
$\theta\in[r-Cr^2,r]$. Decrease $r_0$ so that $Cr_0<1$.
First subtract repeated copies of one fixed positive jump smaller than
$r_0$ from a representative $y\in[0,2\pi)$, leaving a nonnegative
residual below $r_0$. This takes a bounded number of copies independent
of $y$. At each subsequent step subtract the jump given above.
If $r_j$ is the remaining distance after $j$ of these steps, then
$Cr_j\le(Cr_0)^{2^j}$. Since $\log(1/a)=L+O(\log L)$,
every $y$ is within $a$ of a sum of
\[
 m\le\frac{\log L}{\log 2}+O_\chi(1)
\]
jumps. Exact termination is allowed.

Delete one unit of jump intensity for each distinct angle used. Denote
the number of such angles by $d\le m$, their required repeat counts by
$k_i$, and the remaining density by $g_t$. The change in any Fourier
exponent has modulus at most $2d$. From the central Fourier bound
$\sum_n e^{-t\Re\lambda(n)}\ll 1/a$ we obtain, uniformly in all
choices,
\[
 \|g_t-p_t\|_\infty\ll a^{-1}\,td\,e^{2td}.
\]
Thus $g_t(z)\gg a^{-1}$ for $|z|\le a$, since the scaled central
density converges uniformly there to a strictly positive Cauchy
profile. The deleted atoms independently execute the required counts
with probability $e^{-td}t^m/\prod_i k_i!\ge e^{-tm}t^m/m!$.
The contribution of this event to the full kernel proves
\eqref{eq:uniformfloor}, using $p_t(0)\asymp a^{-1}$.

\paragraph{Metric estimates for finite sums.}
Recall that $S_j$ consists of sums of at most $j$ jump locations modulo
$2\pi$, including zero. Write $F_j(\delta)$ for the length of its
closed $\delta$-neighbourhood. The two-sided zero count gives
\[
 \nu\{|u|>r\}\ll r^{-1}\log(e/r),\qquad
 V_r:=\int_{|u|\le r}u^2\,d\nu(u)\ll r\log(e/r)
 \qquad(0<r\le1).
\]
The second inequality follows from the first by integration of the
tail. If every summand has norm at most $r$, the sum is within $jr$
of zero. Otherwise remove a summand of norm greater than $r$.
Consequently
\[
 F_j(\delta)\le 2(jr+\delta)
       +C r^{-1}\log(e/r)F_{j-1}(\delta),\qquad F_0(\delta)=2\delta.
\]
Fix $0<\beta,\varepsilon<1$ and let
$\delta=e^{-\beta n}$, $k=\lfloor(1-\varepsilon)\log n/\log 2\rfloor$.
At step $j$ use $r=\delta^{1/2^j}$. Induction in the last display yields
$F_j(\delta)\le 2(C'(k+1)(n+1))^j\delta^{1/2^j}$ for $0\le j\le k$.
In particular $F_k(\delta)\le\exp(-\beta n^\varepsilon/2)$ for large
$n$, which is summable. By Borel--Cantelli, for almost every $y$ and
all sufficiently large integers $n$, the distance from $y$ to $S_k$
exceeds $\delta$.

\paragraph{A density bound uniform in the time slab.}
For such a $y$ take any $e^{-n-1}\le t\le e^{-n}$.
Truncate the lifted jumps at $\eta=\delta/(16k)$ and put
$\Lambda_\eta=\nu\{|u|>\eta\}$. Then
$t\Lambda_\eta\le\exp(-(1-\beta)n+O(\log n))$.
With all jumps compensated, the small-jump drift is
$b_\eta=b-\int_{|u|>\eta}u\,d\nu$; absolute summation on
$|u|>\eta$ gives $|b_\eta|\ll1+\log^2(e/\eta)$.
Set $s=(1-\beta)n/(2\eta)$. For the small-jump lift $Y_v$ and
$v\le t$, the exponential moment is bounded by
\[
 \mathbb E e^{\pm sY_v}
 \le \exp\bigl(t s|b_\eta|+t s^2 e^{s\eta}V_\eta/2\bigr)\le 2
\]
for large $n$. Indeed the first term in the exponent is exponentially
small with rate $1-\beta$, and the second has rate $(1-\beta)/2$, up
to powers of $n$. Hence
$\mathbb P(|Y_{t/2}|\ge\delta/2)\le4e^{-s\delta/2}$.
If $g_v$ denotes its projected density then
$\|g_{t/2}\|_\infty\ll 1/a(t)$ by Fourier summation, since
$t\Lambda_\eta=o(1)$ and $a(t/2)\asymp a(t)$.
Splitting the convolution $g_t=g_{t/2}*g_{t/2}$ according as one
of its two arguments has circle distance at least $\delta/2$ from
zero gives
\[
 \sup_{\|z\|_{\mathbb T}\ge\delta}g_t(z)
 \ll a(t)^{-1} e^{-s\delta/2}.
\]
Every sum of at most $k$ large jumps lies in $S_k$, so this bounds
the contributions at $y$ from that entire part of the Poisson mixture.
The complementary part is bounded by $O(1/a(t))$ times the probability
of at least $k+1$ large jumps. Therefore
\[
 a(t)p_t(y)\ll e^{-4(1-\beta)nk}
   +(t\Lambda_\eta)^{k+1}.
\]
Divide negative logarithms by $L\log L$, use $n\le L\le n+1$ and
$p_t(0)\asymp 1/a(t)$. The lower limit is at least
$(1-\beta)(1-\varepsilon)/\log 2$. Sending both parameters to zero
through a countable sequence and applying \eqref{eq:uniformfloor}
proves \eqref{eq:typicallaw}.

\begin{corollary}[How many zero angles approximate a prescribed phase?]
Let $D_y(\delta)=\min\{j:\dist(y,S_j)\le\delta\}$. Then
\[
 \sup_{y\in\mathbb T}D_y(\delta)\le
 \frac{\log\log(1/\delta)}{\log2}+O_\chi(1).
\]
For Lebesgue-almost every $y$, and also with $D_y$ replaced by its
supremum over the circle,
\[
 \lim_{\delta\downarrow0}
 \frac{D_y(\delta)}{\log\log(1/\delta)}=\frac1{\log2}.
\]
\end{corollary}
\begin{proof}
The greedy construction above reaches the required tolerance after the
displayed number of jumps. For the reverse bound apply the covering
estimate with $\delta=e^{-n/2}$ and
$k=\lfloor(1-\varepsilon)\log n/\log 2\rfloor$.
Almost every $y$ has $D_y(e^{-n/2})>k$ eventually. Monotonicity
interpolates the intervening tolerances, and a countable sequence of
$\varepsilon\downarrow0$ gives the claim. Taking any such $y$ gives
the lower bound for the supremum.
\end{proof}

\begin{corollary}[A simultaneous bound on the ordinates]\label{heightcap}
The uniform approximation bound can be attained using only zeros with
$|\gamma|\le C_\chi\delta^{-1/2}$. The same almost-everywhere and worst-case
leading terms hold under this height restriction, for this choice of $C_\chi$.
\end{corollary}
\begin{proof}
Use \eqref{rv} for both signs. For all small $r$ the positive angles, and
also the magnitudes of negative angles, meet $[r-Cr^2,r]$ with a common
constant $C$. Choose $u=c\sqrt\delta$, with $9Cc^2<1$. Run the greedy
algorithm until the residual $r$ belongs to $[0,u]$. After shrinking
$r_0$ so $Cr_0\le1/2$, each step from a residual above $u$ uses an angle
at least $u/2$. The number of steps is at most
$\log_2\log(1/u)+O_\chi(1)=\log_2\log(1/\delta)+O_\chi(1)$.
To finish, take a positive angle $b\in[3u-9Cu^2,3u]\subset[2u,3u]$.
Choose a negative angle $-v$ whose magnitude belongs to
$[b-r-C(b-r)^2,b-r]$. Then $|b-v-r|\le9Cu^2<\delta$ and $v\ge u/2$
for small $\delta$. These last two jumps allow a tiny residual to be
approximated without using a correspondingly tiny angle. All angles used,
including the fixed preliminary step, have norm at least $u/2$. Formula
\eqref{cayley} gives $|\gamma|\le1/|\theta_\gamma|\le2/u$. The lower
bound for the unrestricted number of factors still applies.
\end{proof}

In particular, to approximate a specified point of the unit circle within angular error $10^{-D}$, one needs at most $\log_2 D+O_\chi(1)$ factors $1-1/\rho$ from the zeros of $L(s,\chi)$, allowing repetition. For almost every target the least number divided by $\log_2 D$ tends to one. The upper bound comes from approximating the residual to its \emph{square} at each step. The lower bound counts all sums, so a non-greedy choice of zeros cannot improve the leading constant on a positive-measure set. The estimates for angles made from real ordinates use only \eqref{rv}. The identification with the actual factors $1-1/\rho$ uses Dirichlet GRH.

\section{Thirteen asymmetric values for a quartic character}\label{example}
Let $\chi$ be the primitive quartic character modulo five with $\chi(2)=\I$. On $1,2,3,4$, the conjugate character has values $1,-\I,\I,-1$. The interval programs accompanying this note prove that $L(s,\bar\chi)$ has exactly one zero throughout the critical strip with $|\gamma|\le4.14$. It is simple and satisfies
\[
    4.13290370521275 <\gamma_*<4.13290370521295 .
\]
Hence its angle $y_*=\theta_{\gamma_*}$ is the unique jump of largest absolute size for $\bar\chi$, with multiplicity one. The program also proves
\[
 0.24079042878934<y_*<0.24079042878938,\qquad 13y_*<\pi<14y_*.
\]
For $1\le k\le13$, a sum of at most $k$ jumps has absolute value at most $ky_*<\pi$ before reduction to the circle. Equality at the positive endpoint requires exactly $k$ copies of $+y_*$. The spectrum for $\chi$ has its unique extremal jump at $-y_*$ and cannot attain $+ky_*$ with at most $k$ jumps. Thus
\begin{equation}\label{quartic}
 \frac{p_{\bar\chi,t}(ky_*)}{p_{\bar\chi,t}(0)}
       =\frac{t^k}{k!}+O_k(t^{k+1}),\qquad
 \frac{p_{\chi,t}(ky_*)}{p_{\chi,t}(0)}=O_k(t^{k+1}).
\end{equation}
The conclusions exchange signs at $-ky_*$. For example, at the second point the first coefficient is $1/2$, and at the thirteenth it is $1/6227020800$. The denominators at zero agree for the two conjugate characters.

\paragraph{The finite certificates.}
The nearby ProveIt low-zero report \cite{lowzeros} supplied a starting location for the quartic zero. The present checks cover both signs all the way to $4.14$, as required for the extremal argument. Let $c=1/2+4.13290370521285\I$. The first program evaluates the four residue-class series by Euler--Maclaurin through $B_{2j}$ with $n$ initial terms per class. For $\Re s\ge1/2$, $|s|\le5$ and $m=2j$, the sum of the four remainder moduli is bounded by
\[
 \frac{4|B_m|\,5^{m-1}(5\cdot6\cdots(4+m))}{m!}
 \frac{(5n)^{1/2-m}}{m-1/2}.
\]
This follows from the integral remainder, $|(s)_m|\le(5)_m$ and the Fourier-series bound $\sup_x|B_m(\{x\})|\le|B_m|$. All inequalities used for a decision are evaluated with interval arithmetic from rational endpoints. Periodicity gives $|A(x)|=|\sum_{n\le x}\bar\chi(n)|\le\sqrt2$. Differentiating $s\int_1^\infty A(x)x^{-s-1}\dd x$ bounds $|L'|$ by $28$ on the tested line segment and $|L''|$ by $112$ in the larger disk and on the forward-difference segment at $c$.

With $n=32$, $j=6$ the code encloses $|L(c)|<5\cdot10^{-15}$ and $|L'(c)|>1$. The derivative enclosure corrects a real forward difference of step $10^{-5}$ by the bound for $L''$. Rouch\'e's theorem on circles of radii $10^{-13}$ and $10^{-3}$ gives one simple zero in each disk, hence the same zero. Symmetry under $s\mapsto1-\bar s$ places it on the critical line. Outside the larger disk the first program covers the remaining line segments by 301 closed intervals. On each interval of radius $r$ it proves $|L(\text{midpoint})|>28r$ with $n=12$, $j=4$.

The second program traverses the boundary of $[-0.01,1.01]+\I[-4.14,4.14]$. Its interval enclosures place the image of each accepted boundary segment in an open half-plane avoiding zero. Each image segment can therefore be replaced by its endpoint chord without changing the winding number. The 1,918 accepted segments give winding number one by signed crossings of the positive real ray. As $L(s,\bar\chi)$ is entire, this certifies the local strip count independently of GRH. The archive contains both programs and their recorded outputs.

\section{Prior results and further questions}
Lagarias \cite{lagarias} established the complex Li coefficients, their asymptotics and the discrete Fourier support of an interpolating function under GRH. The source article \cite{article} supplies the symmetric central profile, including its expansion to every logarithmic order. Picard \cite{picard,picard2} related small-time density to finite-jump accessibility for more general processes and obtained orders of decay. Hasebe and Huang \cite{huang} analyze nonuniqueness of circle L\'evy measures, including the two opposite Poisson jumps above, attributed there to C\'ebron. Chen, Jiang and Wu \cite{chen} give the fixed-depth Minkowski dimension $1-2^{-j}$ for $j$-fold sums of the harmonic sequence. Their result provides a geometric precedent for the successive squaring scale. Here the sets are signed Cayley angles of zeros, the depth grows as the tolerance shrinks, and the covering argument and the jump probabilities together give the sharp heat-density constant.

The specific contributions are the one-time identification and its phase ambiguity, the multiplicity-weighted pointwise profiles, and the almost-everywhere and minimum laws with their common coefficient. We found no identical result in the sources examined. Global priority would require further literature review. The proofs settle inverse and approximation questions about the Li heat kernel. Their conclusions concern this specific arithmetic model.

\begin{enumerate}[leftmargin=*,label=\arabic*.]
\item Can the one-time theorem yield practical bounds on the number and precision of Fourier modes sufficient for characters of conductor at most $Q$?
\item After subtracting $L\log L/\log2$ from the generic negative log-density, what is the next order? What is the Hausdorff dimension of the exceptional set for a prescribed decay rate?
\item Beyond the extremal construction modulo five, which equal products of Cayley factors from different zero tuples can be proved? The multiplicity formula adds their ordered weights without assumptions on such relations.
\end{enumerate}

\small
\begin{thebibliography}{10}
\bibitem{lean} L. Sernau and collaborators. \emph{GRH for Dirichlet $L$-functions, in Lean 4}, rewriting of the proof of L. Alp\"oge developed with Claude. Commit \texttt{72fead102f66}, 25 Nov.\ 2026.
\href{\leanbase/Statements.lean}{Public statements} and \href{\leanbase/Main.lean}{proved declarations}.
\bibitem{mit} A. Sutherland. \emph{Verification files for the shorter Lean proof of RH and of GRH for Dirichlet $L$-functions}, 27 Nov.\ 2026. \url{https://math.mit.edu/~drew/shortgrh/}.
\bibitem{article} V. Reshetnikov, ProveIt collection. \emph{Beyond the Riemann hypothesis for Dirichlet $L$-functions: consequences, transfers and limits}. Commit \texttt{1a50d416ba30}.
\href{\artbase/Article/README.md}{Article overview},
\href{\artbase/Article/sections/src.tex}{source theorem} and
\href{\artbase/Article/sections/li.tex}{Li coefficients, subsection on the Cayley angles}.
\bibitem{lowzeros} ProveIt research collection at the same commit.
\href{\artbase/Research/dirichlet_low_zero_identification/character_zeros.tex}{\emph{Distinguishing Dirichlet characters from low zeros}}, including the quartic example.
\href{\artbase/Research/dirichlet_zeros_and_abelian_fields/zero_windows.tex}{\emph{Distinguishing Dirichlet characters and abelian number fields by low zeros}}.
\bibitem{lagarias} J. C. Lagarias. Li coefficients for automorphic $L$-functions. \emph{Ann. Inst. Fourier} \textbf{57} (2007), 1689--1740. \href{https://doi.org/10.5802/aif.2311}{doi:10.5802/aif.2311}.
\bibitem{hasebe} T. Hasebe and I. Hotta. Additive processes on the unit circle and Loewner chains. \emph{Int. Math. Res. Not.} (2022), \S4.1. \href{https://doi.org/10.1093/imrn/rnab157}{doi:10.1093/imrn/rnab157}.
\bibitem{huang} T. Hasebe and H.-W. Huang. Limit theorems and wrapping transforms in bi-free probability theory. \href{https://arxiv.org/abs/2007.02775v1}{arXiv:2007.02775v1}, \S8.
\bibitem{picard} J. Picard. Density in small time for L\'evy processes. \emph{ESAIM Probab. Stat.} \textbf{1} (1997), 357--389. \href{https://www.numdam.org/item/PS_1997__1__357_0/}{Article}.
\bibitem{picard2} J. Picard. Density in small time at accessible points for jump processes. \emph{Stochastic Process. Appl.} \textbf{67} (1997), 251--279. \href{https://lmbp.uca.fr/~picard/small_time.pdf}{Author's copy}.
\bibitem{chen} H. Chen, L. Jiang and Y. Wu. Representations of rational numbers and Minkowski dimension. \href{https://arxiv.org/abs/2510.17112}{arXiv:2510.17112}, Theorem 1.4.
\end{thebibliography}
\end{document}
